% ============================================================
% AULA EM BEAMER — Números Complexos
% Aula online expositiva com espaços para desenvolvimento manual
% Base visual adaptada de: template_beamer_paleta_montserrat_ubuntu
% ============================================================

\documentclass[aspectratio=169]{beamer}

% ------------------------------------------------------------
% 1. Pacotes básicos
% ------------------------------------------------------------
\usepackage[T1]{fontenc}
\usepackage[brazilian]{babel}
\usepackage{amsmath,amssymb}
\usepackage{graphicx}
\usepackage{tikz}
\usepackage{xcolor}
\usepackage{booktabs}
\usepackage{array}
\usepackage{multicol}
\usepackage{url}
\usepackage{iftex}
\usepackage{standalone}

% ------------------------------------------------------------
% 2. Tema visual base
% ------------------------------------------------------------
\usetheme{Madrid}

% Fontes do template
% Recomendação: compilar com XeLaTeX ou LuaLaTeX para melhor uso de fontes.
\ifPDFTeX
  \usepackage[defaultfam,tabular,lining]{montserrat}
  \renewcommand{\familydefault}{\sfdefault}
  \newcommand{\FonteDestaque}{\sffamily\mdseries}
\else
  \usepackage{fontspec}
  \IfFontExistsTF{Montserrat}{%
    \setmainfont{Montserrat}
    \setsansfont{Montserrat}
  }{%
    \IfFontExistsTF{Noto Sans}{%
      \setmainfont{Noto Sans}
      \setsansfont{Noto Sans}
    }{%
      \setmainfont{Latin Modern Sans}
      \setsansfont{Latin Modern Sans}
    }
  }
  \IfFontExistsTF{Ubuntu}{%
    \newfontfamily\UbuntuFont{Ubuntu}
  }{%
    \IfFontExistsTF{Noto Sans}{%
      \newfontfamily\UbuntuFont{Noto Sans}
    }{%
      \newfontfamily\UbuntuFont{Latin Modern Sans}
    }
  }
  \newcommand{\FonteDestaque}{\UbuntuFont\mdseries}
\fi

\usefonttheme{professionalfonts}
\setbeamerfont{normal text}{family=\sffamily}
\setbeamerfont{title}{family=\FonteDestaque,series=\mdseries,size=\LARGE}
\setbeamerfont{subtitle}{family=\FonteDestaque,series=\mdseries,size=\large}
\setbeamerfont{frametitle}{family=\FonteDestaque,series=\mdseries,size=\large}
\setbeamerfont{block title}{family=\FonteDestaque,series=\mdseries}
\setbeamerfont{enumerate item}{series=\bfseries}
\setbeamerfont{enumerate subitem}{series=\bfseries}

% ------------------------------------------------------------
% 3. Paleta visual
% ------------------------------------------------------------
\definecolor{c_red}{HTML}{e63616}
\definecolor{c_yellow}{HTML}{fbbc1e}
\definecolor{c_white}{HTML}{f2f2f2}
\definecolor{c_blue1}{HTML}{2158a5}
\definecolor{c_blue2}{HTML}{152f4d}
\definecolor{c_blue3}{HTML}{0f1626}

\colorlet{ThemePrimary}{c_blue2}
\colorlet{ThemeSecondary}{c_blue1}
\colorlet{ThemeDark}{c_blue3}
\colorlet{ThemeAccent}{c_yellow}
\colorlet{ThemeAlert}{c_red}
\colorlet{ThemeLight}{c_white}

\setbeamercolor{normal text}{fg=ThemeDark,bg=ThemeLight}
\setbeamercolor{background canvas}{bg=ThemeLight}
\setbeamercolor{palette primary}{bg=ThemePrimary, fg=ThemeLight}
\setbeamercolor{palette secondary}{bg=ThemeSecondary, fg=ThemeLight}
\setbeamercolor{palette tertiary}{bg=ThemeDark, fg=ThemeLight}
\setbeamercolor{palette quaternary}{bg=ThemeAlert, fg=ThemeLight}
\setbeamercolor{frametitle}{bg=ThemePrimary, fg=ThemeLight}
\setbeamercolor{title}{fg=ThemeLight}
\setbeamercolor{subtitle}{fg=ThemeAccent}
\setbeamercolor{author}{fg=ThemeLight}
\setbeamercolor{institute}{fg=ThemeLight}
\setbeamercolor{date}{fg=ThemeLight}
\setbeamercolor{structure}{fg=ThemeSecondary}
\setbeamercolor{alerted text}{fg=ThemeAlert}
\setbeamercolor{block title}{bg=ThemeSecondary, fg=ThemeLight}
\setbeamercolor{block body}{bg=ThemeLight, fg=ThemeDark}
\setbeamercolor{itemize item}{fg=ThemeAccent}
\setbeamercolor{itemize subitem}{fg=ThemeSecondary}
\setbeamercolor{section in toc}{fg=ThemePrimary}

\setbeamertemplate{navigation symbols}{}

% ------------------------------------------------------------
% 4. Logo opcional no canto superior direito
% ------------------------------------------------------------
\newcommand{\logounb}{logoUnB.png}

\addtobeamertemplate{frametitle}{}{%
    \begin{tikzpicture}[remember picture,overlay]
        \node[anchor=north east, xshift=-0.25cm, yshift=-0.15cm]
            at (current page.north east)
            {\IfFileExists{\logounb}{\includegraphics[height=0.65cm]{\logounb}}{}};
    \end{tikzpicture}
}

% ------------------------------------------------------------
% 5. Comandos auxiliares
% ------------------------------------------------------------
\graphicspath{{./}}
\newcommand{\sen}{\operatorname{sen}}

% Área realmente vazia para escrita manual em mesa digitalizadora/tablet.
% Use dentro de slides de exemplo, demonstração ou resolução.
\newcommand{\EspacoEscrita}[1]{%
  \vspace{0.25em}
  \begin{minipage}[t][#1][t]{\linewidth}
  \end{minipage}
}

% Etiquetas discretas de tempo.
\newcommand{\Tempo}[1]{\hfill {\scriptsize\color{ThemeAccent}\textbf{#1}}}

% Itens do sumário sanfonado por seção.
\newcommand{\secaoPassada}[1]{{\color{ThemeSecondary}\checkmark\ #1}}
\newcommand{\secaoAtual}[1]{{\color{ThemeAccent}\textbf{\large #1}}}
\newcommand{\secaoFutura}[1]{{\color{ThemeDark!55}#1}}

% Slide de sumário local: o bloco da seção ocupa cerca de 35% da largura.
\newcommand{\MapaSecao}[3]{%
\begin{frame}{Roteiro da aula — seção atual}
\vspace{-0.55em}
\begin{columns}[T,onlytextwidth]
\column{0.54\textwidth}
\footnotesize
\textbf{Percurso geral}\par\vspace{0.25em}
#1
\column{0.03\textwidth}
\column{0.35\textwidth}
\footnotesize
\textbf{Itens da seção}\par\vspace{0.25cm}
\begin{block}{#2}
#3
\end{block}
\column{0.03\textwidth}
\end{columns}
\end{frame}
}

% ------------------------------------------------------------
% 6. Metadados da apresentação
% ------------------------------------------------------------
\title[Números Complexos]{Números Complexos}
\subtitle{Conceito, operações e representação geométrica}
\author[Carlos Eduardo da Silva Papa]{Prof. \textit{Cadu}}
% \institute[Ensino Médio]{Aula online expositiva}

% ============================================================
% INÍCIO DO DOCUMENTO
% ============================================================
\begin{document}

% ------------------------------------------------------------
% CAPA
% ------------------------------------------------------------
\begin{frame}
  \titlepage
\end{frame}

% ------------------------------------------------------------
% SUMÁRIO PRINCIPAL
% ------------------------------------------------------------
\begin{frame}{Sumário principal}
\vspace{-0.65em}
\begin{columns}[T,onlytextwidth]

\column{0.49\textwidth}
\small
\textbf{\color{ThemeDark}I — Construção do conjunto complexo}
\begin{itemize}
  \item O ``por que`` dos números complexos
  \item Forma algébrica
  \item Igualdade e classificação
\end{itemize}

\vspace{0.35em}
\textbf{\color{ThemeDark}II — Cálculo algébrico}
\begin{itemize}
  \item Adição, subtração e multiplicação
  \item Potências de $i$
\end{itemize}

\column{0.49\textwidth}
\small
\textbf{\color{ThemeDark}III — Geometria e novas formas}
\begin{itemize}
  \item Plano complexo, módulo e argumento
  \item Conjugado e divisão
  \item Representação polar e operações
\end{itemize}

\vspace{0.35em}
\textbf{\color{ThemeDark}IV — Aplicação e encerramento}
\begin{itemize}
  \item Equação com raízes complexas
  \item Revisão e dever de casa
\end{itemize}

\vspace{0.7em}
% \begin{block}{Duração total}
% \centering\textbf{100 minutos}
% \end{block}

\end{columns}
\end{frame}

% ============================================================
\section{O ``por que`` dos números complexos}
% ============================================================

\MapaSecao{
\begin{itemize}
  \item \secaoAtual{O ``por que`` dos números complexos}
  \item \secaoFutura{Forma algébrica}
  \item \secaoFutura{Igualdade e classificação}
  \item \secaoFutura{Operações algébricas}
  \item \secaoFutura{Plano complexo}
  \item \secaoFutura{Conjugado e divisão}
  \item \secaoFutura{Forma polar}
  \item \secaoFutura{Aplicação}
  \item \secaoFutura{Fechamento}
\end{itemize}
}{Por que os complexos?}{
\begin{itemize}
  \item Limitação do conjunto dos reais.
  \item Definição da unidade imaginária.
  \item Lugar de $\mathbb{C}$ entre os conjuntos numéricos.
\end{itemize}
}

\begin{frame}{Resolução em aula: por que ampliar os números reais?}
\vspace{-0.2em}
\small
Resolver e discutir no quadro digital:
\[
  x^2 + 1 = 0
\]
\EspacoEscrita{0.58\textheight}
\end{frame}

\begin{frame}{Unidade imaginária}
\vspace{-0.4em}
\small
Define-se a unidade imaginária por:

\vspace{0.8em}
\begin{block}{Definição}
\[
  i^2=-1
\]
\end{block}

\vspace{0.5em}
\begin{block}{Consequência imediata}
\[
  \sqrt{-1}=i
\]
\end{block}
\end{frame}

\begin{frame}{Os complexos entre os conjuntos numéricos}
\centering
\resizebox{0.48\textwidth}{!}{%
\begin{tikzpicture}[font=\Large]
% Bounding box FIXO: garante que o tamanho do desenho nunca muda
% entre os passos do overlay (mesmo com menos elementos visíveis).
\useasboundingbox (-2.5,-2.5) rectangle (4,2.5);

% Quadrado dos Complexos — aparece por último
\fill<6->[green!25] (-2.5,-2.5) rectangle (4,2.5);
\draw<6->[very thick] (-2.5,-2.5) rectangle (4,2.5);
\node<6-> at (-2,2) {$\mathbb{C}$};

% Quadrado dos Reais
\fill<4->[green!5] (-2,-2) rectangle (3.5, 1.5);
\draw<4->[very thick] (-2,-2) rectangle (3.5, 1.5);
\node<4-> at (-1.5,1) {$\mathbb{R}$};

% Círculos internos (Racionais, Inteiros, Naturais)
\fill<3->[yellow!20] (-0.6,-0.5) circle (1.2);
\fill<2->[orange!20] (-0.6,-0.9) circle (0.8);
\fill<1->[red!20] (-0.6,-1.3) circle (0.4);
\draw<3->[very thick] (-0.6,-0.5) circle (1.2);
\draw<2->[very thick] (-0.6,-0.9) circle (0.8);
\draw<1->[very thick] (-0.6,-1.3) circle (0.4);
\node<3-> at (-0.6, 0.3) {$\mathbb{Q}$};
\node<2-> at (-0.6,-0.5) {$\mathbb{Z}$};
\node<1-> at (-0.6,-1.3) {$\mathbb{N}$};

% Círculo dos Irracionais dentro do retângulo dos reais
\fill<5->[purple!20] (2,-0.5) circle (1);
\draw<5->[very thick] (2,-0.5) circle (1);
\node<5-> at (2,-0.5) {$\mathbb{I}$};
\end{tikzpicture}%
}

\vspace{-0.2em}
\small
\[
  \mathbb{N}
  \visible<2->{{}\subset\mathbb{Z}}
  \visible<3->{{}\subset\mathbb{Q}}
  \visible<4->{{}\subset\mathbb{R}}
  \visible<6->{{}\subset\mathbb{C}}
\]
\visible<5->{\footnotesize (com $\mathbb{I}\subset\mathbb{R}$, os irracionais)}
\end{frame}

% ============================================================
\section{Forma algébrica}
% ============================================================

\MapaSecao{
\begin{itemize}
  \item \secaoPassada{O ``por que`` dos números complexos}
  \item \secaoAtual{Forma algébrica}
  \item \secaoFutura{Igualdade e classificação}
  \item \secaoFutura{Operações algébricas}
  \item \secaoFutura{Plano complexo}
  \item \secaoFutura{Conjugado e divisão}
  \item \secaoFutura{Forma polar}
  \item \secaoFutura{Aplicação}
  \item \secaoFutura{Fechamento}
\end{itemize}
}{Forma algébrica}{
\begin{itemize}
  \item Definição geral: $z=a+bi$.
  \item Parte real e parte imaginária.
  \item Identificação em exemplos.
\end{itemize}
}

\begin{frame}{Forma algébrica de um número complexo}
\vspace{-0.3em}
\small
Todo número complexo pode ser escrito na forma:

\vspace{0.5em}
\begin{block}{Forma algébrica}
\[
  z=a+bi
\]
\end{block}

\vspace{0.4em}
\begin{columns}[T,onlytextwidth]
\column{0.48\textwidth}
\begin{block}{Parte real}
\[
  \operatorname{Re}(z)=a
\]
\end{block}
\column{0.48\textwidth}
\begin{block}{Parte imaginária}
\[
  \operatorname{Im}(z)=b
\]
\end{block}
\end{columns}
\end{frame}

\begin{frame}{Resolução em aula: identificar componentes}
\vspace{-0.2em}
\small
Para cada número, identificar $a$, $b$, $\operatorname{Re}(z)$ e $\operatorname{Im}(z)$:
\[
  z_1=4-3i, \qquad z_2=-8+5i, \qquad z_3=-9i, \qquad z_4=11
\]
\begin{block}{Atenção}
A parte imaginária é o coeficiente de $i$, não o termo completo com $i$.
\end{block}
\EspacoEscrita{0.39\textheight}
\end{frame}

% ============================================================
\section{Igualdade e classificação}
% ============================================================

\MapaSecao{
\begin{itemize}
  \item \secaoPassada{O ``por que`` dos números complexos}
  \item \secaoPassada{Forma algébrica}
  \item \secaoAtual{Igualdade e classificação}
  \item \secaoFutura{Operações algébricas}
  \item \secaoFutura{Plano complexo}
  \item \secaoFutura{Conjugado e divisão}
  \item \secaoFutura{Forma polar}
  \item \secaoFutura{Aplicação}
  \item \secaoFutura{Fechamento}
\end{itemize}
}{Igualdade e classificação}{
\begin{itemize}
  \item Critério de igualdade entre complexos.
  \item Número real puro, imaginário puro e complexo não real.
\end{itemize}
}

\begin{frame}{Igualdade entre números complexos}
\vspace{-0.3em}
\small
Dois números complexos são iguais quando suas partes reais e imaginárias são iguais.

\vspace{0.6em}
\begin{block}{Critério de igualdade}
\[
  a+bi=c+di
\]
se, e somente se,
\[
  a=c \quad \text{e} \quad b=d.
\]
\end{block}
\end{frame}

\begin{frame}{Classificação dos números complexos}
\vspace{-0.4em}
\small
\begin{columns}[T,onlytextwidth]
\column{0.32\textwidth}
\begin{block}{Real puro}
\[
  z=a
\]
\centering $b=0$
\end{block}

\column{0.32\textwidth}
\begin{block}{Imaginário puro}
\[
  z=bi
\]
\centering $a=0$ e $b\neq 0$
\end{block}

\column{0.32\textwidth}
\begin{block}{Complexo não real}
\[
  z=a+bi
\]
\centering $b\neq 0$
\end{block}
\end{columns}
\end{frame}

\begin{frame}{Resolução em aula: igualdade e classificação}
\vspace{-0.2em}
\small
Determinar $x$ e $y$:
\[
  (2x-1)+(y+3)i=7-4i
\]

\vspace{0.5em}
Classificar os números:
\[
  -6, \qquad 4i, \qquad 5-2i, \qquad 0, \qquad -3i.
\]
\EspacoEscrita{0.46\textheight}
\end{frame}

% ============================================================
\section{Operações algébricas}
% ============================================================

\MapaSecao{
\begin{itemize}
  \item \secaoPassada{O ``por que`` dos números complexos}
  \item \secaoPassada{Forma algébrica}
  \item \secaoPassada{Igualdade e classificação}
  \item \secaoAtual{Operações algébricas}
  \item \secaoFutura{Plano complexo}
  \item \secaoFutura{Conjugado e divisão}
  \item \secaoFutura{Forma polar}
  \item \secaoFutura{Aplicação}
  \item \secaoFutura{Fechamento}
\end{itemize}
}{Operações algébricas}{
\begin{itemize}
  \item Adição e subtração.
  \item Multiplicação.
  \item Potências de $i$.
\end{itemize}
}

\begin{frame}{Adição e subtração}
\vspace{-0.45em}
\small
Somam-se ou subtraem-se separadamente as partes reais e imaginárias.

\begin{block}{Regra geral}
\[
  (a+bi)\pm(c+di)=(a\pm c)+(b\pm d)i
\]
\end{block}
\begin{columns}[T,onlytextwidth]
\column{0.47\textwidth}
\textbf{Exemplo 1}
\[(3+2i)+(5-4i)\]
\column{0.47\textwidth}
\textbf{Exemplo 2}
\[(-6+7i)-(2-5i)\]
\end{columns}
\EspacoEscrita{0.25\textheight}
\end{frame}

\begin{frame}{Multiplicação}
\vspace{-0.45em}
\small
Na multiplicação, aplica-se a distributiva e substitui-se $i^2$ por $-1$.

\begin{block}{Regra geral}
\[
  (a+bi)(c+di)=(ac-bd)+(ad+bc)i
\]
\end{block}
\begin{columns}[T,onlytextwidth]
\column{0.47\textwidth}
\textbf{Exemplo 1}
\[(2+3i)(1-i)\]
\column{0.47\textwidth}
\textbf{Exemplo 2}
\[(-4+i)(3+2i)\]
\end{columns}
\EspacoEscrita{0.25\textheight}
\end{frame}

\begin{frame}{Resolução em aula: potências de $i$}
\vspace{-0.8em}
\small
% As potências repetem o ciclo $i,-1,-i,1$, de período $4$. Use o resto da divisão do expoente por $4$ para calcular:
\[
  i^{17}, \qquad i^{38}, \qquad i^{101}, \qquad i^{2026}
\]
\EspacoEscrita{0.50\textheight}
\end{frame}

% ============================================================
\section{Plano complexo}
% ============================================================

\MapaSecao{
\begin{itemize}
  \item \secaoPassada{O ``por que`` dos números complexos}
  \item \secaoPassada{Forma algébrica}
  \item \secaoPassada{Igualdade e classificação}
  \item \secaoPassada{Operações algébricas}
  \item \secaoAtual{Plano complexo}
  \item \secaoFutura{Conjugado e divisão}
  \item \secaoFutura{Forma polar}
  \item \secaoFutura{Aplicação}
  \item \secaoFutura{Fechamento}
\end{itemize}
}{Plano complexo}{
\begin{itemize}
  \item Representação no plano.
  \item Módulo e argumento.
  \item Formas retangular e polar.
\end{itemize}
}

\begin{frame}{Plano de Argand-Gauss}
\vspace{-0.4em}
\small
O número $z=a+bi$ corresponde ao ponto $P=(a,b)$.

\vspace{0.35em}
\begin{columns}[T,onlytextwidth]
\column{0.48\textwidth}
\begin{block}{Eixo horizontal}
Parte real: $a$.
\end{block}
\column{0.48\textwidth}
\begin{block}{Eixo vertical}
Parte imaginária: $b$.
\end{block}
\end{columns}

\vspace{0.35em}
\begin{block}{Correspondência}
\[
  z=a+bi \quad\longleftrightarrow\quad P=(a,b)
\]
\end{block}
\end{frame}

\begin{frame}{Construção em aula: os quatro quadrantes}
\vspace{-0.65em}
\begin{columns}[T,onlytextwidth]
\column{0.34\textwidth}
\small
Representar e identificar o quadrante:
\begin{itemize}
  \item $z_1=3+2i$;
  \item $z_2=-3+2i$;
  \item $z_3=-3-2i$;
  \item $z_4=3-2i$.
\end{itemize}
E os pontos sobre os eixos?

\column{0.62\textwidth}
\centering
\includestandalone[width=0.72\textwidth]{plano_cartesiano_1}
\end{columns}
\end{frame}

\begin{frame}{Módulo e argumento}
\vspace{-0.45em}
\small
Para $z=a+bi\neq0$, o módulo é a distância até a origem e o argumento é o ângulo orientado a partir do eixo real positivo.
\begin{columns}[T,onlytextwidth]
\column{0.48\textwidth}
\begin{block}{Módulo}
\[
  r=|z|=\sqrt{a^2+b^2}
\]
\end{block}
\column{0.48\textwidth}
\begin{block}{Argumento}
\[
  \theta=\arg(z),\qquad \tan\theta=\frac{b}{a}
\]
\end{block}
\end{columns}
\begin{alertblock}{Atenção ao quadrante}
A tangente fornece um ângulo de referência; os sinais de $a$ e $b$ determinam o quadrante correto.
\end{alertblock}
\end{frame}

\begin{frame}{Módulo e argumento nos quatro quadrantes}
\vspace{-0.4em}
\small
\centering
\renewcommand{\arraystretch}{1.35}
\begin{tabular}{cccc}
\toprule
Quadrante & Número & Módulo & Um argumento \\
\midrule
I   & $1+i$  & $\sqrt2$ & $45^\circ$ \\
II  & $-1+i$ & $\sqrt2$ & $135^\circ$ \\
III & $-1-i$ & $\sqrt2$ & $225^\circ$ \\
IV  & $1-i$  & $\sqrt2$ & $315^\circ$ \\
\bottomrule
\end{tabular}

\vspace{0.55em}
\begin{block}{Argumentos equivalentes}
\[
  \theta_k=\theta+360^\circ k,\qquad k\in\mathbb{Z}.
\]
$765^\circ=45^\circ+2\cdot360^\circ$ representa a mesma direção de $45^\circ$.
\end{block}
\end{frame}

\begin{frame}{Construção ampliada: módulo, argumento e voltas}
\vspace{-0.65em}
\begin{columns}[T,onlytextwidth]
\column{0.35\textwidth}
\small
No quadro:
\begin{enumerate}
  \item representar $z=-4+4i$;
  \item traçar o segmento de módulo $|z|$;
  \item marcar $135^\circ$;
  \item indicar também $495^\circ$ e $-225^\circ$.
\end{enumerate}

\column{0.61\textwidth}
\centering
\includestandalone[width=0.70\textwidth]{plano_cartesiano_2}
\end{columns}
\end{frame}

\begin{frame}{Da forma retangular à forma polar}
\vspace{-1cm}
\small
Uma volta completa mede $360^\circ=2\pi$ radianos; portanto,
\[
  180^\circ=\pi\text{ rad},\qquad 90^\circ=\frac{\pi}{2}\text{ rad}.
\]
\begin{columns}[T,onlytextwidth]
\column{0.45\textwidth}
\begin{block}{Forma retangular}

\[
  z=a+bi
\]
\end{block}
\column{0.51\textwidth}
\begin{block}{Formas polares}
\[
  z=r(\cos\theta+i\sen\theta)
\]
\[
  z=r\,e^{i\theta}
\]
\end{block}
\end{columns}
\centering
\vspace{1cm}
$a=r\cos\theta$ e $b=r\sen\theta$.
\end{frame}

\begin{frame}{Resolução em aula: conversão entre representações}
\vspace{-0.25em}
\small
Converter para a forma polar:
\[
  z_1=1+i,\qquad z_2=-\sqrt3+i,\qquad z_3=-2i.
\]
Converter para a forma retangular:
\[
  z_4=4\left(\cos\frac{\pi}{3}+i\sen\frac{\pi}{3}\right),
  \qquad z_5=3e^{i\pi}.
\]
\EspacoEscrita{0.43\textheight}
\end{frame}

% ============================================================
\section{Conjugado e divisão}
% ============================================================

\MapaSecao{
\begin{itemize}
  \item \secaoPassada{O ``por que`` dos números complexos}
  \item \secaoPassada{Forma algébrica}
  \item \secaoPassada{Igualdade e classificação}
  \item \secaoPassada{Operações algébricas}
  \item \secaoPassada{Plano complexo}
  \item \secaoAtual{Conjugado e divisão}
  \item \secaoFutura{Forma polar}
  \item \secaoFutura{Aplicação}
  \item \secaoFutura{Fechamento}
\end{itemize}
}{Conjugado e divisão}{
\begin{itemize}
  \item Definição de conjugado.
  \item Interpretação geométrica.
  \item Divisão na forma retangular.
\end{itemize}
}

\begin{frame}{Conjugado: definição e propriedade algébrica}
\vspace{-0.85em}
\small
\begin{columns}[T,onlytextwidth]
\column{0.46\textwidth}
\begin{block}{Definição}
Se $z=a+bi$, seu conjugado é
\[
  \overline z=a-bi.
\]
\end{block}

\begin{block}{Exemplo}
\[
  z=-3+2i
  \quad\Longrightarrow\quad
  \overline z=-3-2i.
\]
\end{block}

\column{0.50\textwidth}
\begin{block}{O que muda?}
O sinal da parte imaginária:
\[
  \operatorname{Im}(\overline z)=-\operatorname{Im}(z).
\]
\end{block}

\begin{block}{O que permanece?}
\[
  \operatorname{Re}(\overline z)=\operatorname{Re}(z),
  \qquad |\overline z|=|z|.
\]
\end{block}
\end{columns}

\begin{alertblock}{Produto fundamental}
\[
  z\overline z=(a+bi)(a-bi)=a^2+b^2=|z|^2.
\]
Esse resultado real será usado na divisão de números complexos.
\end{alertblock}
\end{frame}

\begin{frame}{Conjugado no plano complexo}
\vspace{-1cm}
\begin{columns}[T,onlytextwidth]
\column{0.36\textwidth}
\small
Para
\[
  z=-3+2i,
  \qquad
  \overline z=-3-2i,
\]
construir no quadro:
\begin{enumerate}
  \item os pontos $z$ e $\overline z$;
  \item os segmentos até a origem;
  \item o eixo de simetria.
\end{enumerate}

\begin{block}{Leitura geométrica}
$z=(a,b)$ e $\overline z=(a,-b)$ são reflexos em relação ao eixo real.
\end{block}

\column{0.60\textwidth}
\centering
\includestandalone[width=0.82\linewidth]{plano_cartesiano_1}
\end{columns}
\end{frame}

\begin{frame}{Divisão na forma retangular}
\vspace{-0.45em}
\small
Multiplicamos numerador e denominador pelo conjugado do denominador.
\begin{block}{Procedimento}
\[
  \frac{a+bi}{c+di}\cdot\frac{c-di}{c-di}
  =\frac{(a+bi)(c-di)}{c^2+d^2},\qquad c+di\neq0.
\]
\end{block}
\begin{columns}[T,onlytextwidth]
\column{0.47\textwidth}
\textbf{Exemplo 1}
\[\frac{3+2i}{1-i}\]
\column{0.47\textwidth}
\textbf{Exemplo 2}
\[\frac{5-4i}{2+3i}\]
\end{columns}
\EspacoEscrita{0.22\textheight}
\end{frame}

% ============================================================
\section{Operações na forma polar}
% ============================================================

\MapaSecao{
\begin{itemize}
  \item \secaoPassada{O ``por que`` dos números complexos}
  \item \secaoPassada{Forma algébrica}
  \item \secaoPassada{Igualdade e classificação}
  \item \secaoPassada{Operações algébricas}
  \item \secaoPassada{Plano complexo}
  \item \secaoPassada{Conjugado e divisão}
  \item \secaoAtual{Forma polar}
  \item \secaoFutura{Aplicação}
  \item \secaoFutura{Fechamento}
\end{itemize}
}{Operações na forma polar}{
\begin{itemize}
  \item Multiplicação.
  \item Divisão.
  \item Limite para adição e subtração.
\end{itemize}
}

\begin{frame}{Multiplicação e divisão na forma polar}
\vspace{-0.5em}
\small
Se $z_1=r_1e^{i\theta_1}$ e $z_2=r_2e^{i\theta_2}$, então:
\begin{columns}[T,onlytextwidth]
\column{0.48\textwidth}
\begin{block}{Multiplicação}
\[
  z_1z_2=r_1r_2e^{i(\theta_1+\theta_2)}
\]
Multiplicam-se os módulos e somam-se os argumentos.
\end{block}
\column{0.48\textwidth}
\begin{block}{Divisão}
\[
  \frac{z_1}{z_2}=\frac{r_1}{r_2}e^{i(\theta_1-\theta_2)}
\]
Dividem-se os módulos e subtraem-se os argumentos.
\end{block}
\end{columns}
\begin{alertblock}{E a adição ou a subtração?}
Não se somam nem se subtraem módulos e argumentos diretamente. Para $z_1\pm z_2$, primeiro converta para a forma retangular.
\end{alertblock}
\end{frame}

\begin{frame}{Resolução em aula: operações polares}
\vspace{-0.3em}
\small
Considere
\[
  z_1=2e^{i\pi/3}\qquad\text{e}\qquad z_2=3e^{i\pi/6}.
\]
Calcular e apresentar nas formas polar e retangular:
\[
  z_1z_2,\qquad \frac{z_1}{z_2}.
\]
% Para calcular $z_1+z_2$, por que é necessário mudar de representação?
\EspacoEscrita{0.42\textheight}
\end{frame}

% ============================================================
\section{Aplicação}
% ============================================================

\MapaSecao{
\begin{itemize}
  \item \secaoPassada{O ``por que`` dos números complexos}
  \item \secaoPassada{Forma algébrica}
  \item \secaoPassada{Igualdade e classificação}
  \item \secaoPassada{Operações algébricas}
  \item \secaoPassada{Plano complexo}
  \item \secaoPassada{Conjugado e divisão}
  \item \secaoPassada{Forma polar}
  \item \secaoAtual{Aplicação}
  \item \secaoFutura{Fechamento}
\end{itemize}
}{Aplicação}{
\begin{itemize}
  \item Equação do segundo grau.
  \item Discriminante ($\Delta$) negativo.
  \item Raízes complexas.
\end{itemize}
}

\begin{frame}{Aplicação: resolução de uma equação}
\vspace{-0.25em}
\small
Resolver em $\mathbb{C}$:
\[
  x^2-2x+5=0.
\]
\begin{block}{Roteiro}
\begin{enumerate}
  \item calcular o discriminante ($\Delta$);
  \item usar $\sqrt{-1}=i$;
  \item escrever e conferir as duas raízes complexas;
  \item localizar as raízes no plano complexo.
\end{enumerate}
\end{block}
\EspacoEscrita{0.20\textheight}
\end{frame}

\begin{frame}{Aplicação: raízes no plano complexo}
\vspace{-0.65em}
\begin{columns}[T,onlytextwidth]
\column{0.35\textwidth}
\small
Após resolver a equação:
\[
  x_1=1+2i,
  \qquad
  x_2=1-2i.
\]

\begin{block}{Conclusão}
As raízes são números complexos conjugados:
\[
  x_2=\overline{x_1}.
\]
\end{block}

\column{0.61\textwidth}
\vspace{-.25cm}
\centering
\includestandalone[width=0.78\linewidth]{plano_cartesiano_1}
\end{columns}
\end{frame}


\begin{frame}[t]{Números Complexos: Questão de Aplicação}

\vspace{-.5cm}

\footnotesize

\begin{block}{FQEQ(2017)}
No plano cartesiano, a posição de um robô é representada pelo número complexo
$z=x+yi$, em que $x$ indica o deslocamento horizontal e $y$ o deslocamento vertical.

Após receber dois comandos, sua posição final é dada por
\[
z_f=(2+i)(3-2i)+(-1+4i).
\]

Determine a posição final do robô na forma $a+bi$ e calcule sua distância até a origem.
\end{block}

\end{frame}


% ============================================================
\section{Fechamento}
% ============================================================

\MapaSecao{
\begin{itemize}
  \item \secaoPassada{O ``por que´´ dos números complexos}
  \item \secaoPassada{Forma algébrica}
  \item \secaoPassada{Igualdade e classificação}
  \item \secaoPassada{Operações algébricas}
  \item \secaoPassada{Plano complexo}
  \item \secaoPassada{Conjugado e divisão}
  \item \secaoPassada{Forma polar}
  \item \secaoPassada{Aplicação}
  \item \secaoAtual{Fechamento}
\end{itemize}
}{Fechamento}{
\begin{itemize}
  \item Retomada do percurso.
  \item Escolha da representação.
  \item Dever de casa.
\end{itemize}
}

\begin{frame}{Revisão: do problema aos números complexos}
\vspace{-0.45em}
\small
\begin{enumerate}
  \setlength{\itemsep}{5pt}
  \item Equações como $x^2=-1$ exigem ampliar $\mathbb{R}$ para $\mathbb{C}$.
  \item A unidade imaginária satisfaz $i^2=-1$.
  \item Na forma retangular, $z=a+bi$ tem parte real $a$ e parte imaginária $b$.
  \item Igualdade, adição e subtração são feitas comparando as partes correspondentes.
  \item Na multiplicação, usamos a distributiva e substituímos $i^2$ por $-1$.
  \item O conjugado $\overline z=a-bi$ permite retirar $i$ do denominador em divisões.
\end{enumerate}
\end{frame}

\begin{frame}{Revisão: geometria e escolha da representação}
\vspace{-0.45em}
\small
\begin{itemize}
  \setlength{\itemsep}{5pt}
  \item No plano complexo, $z=a+bi$ é o ponto $(a,b)$.
  \item O módulo $|z|=\sqrt{a^2+b^2}$ mede a distância até a origem.
  \item O argumento indica a direção; ângulos que diferem por $360^\circ k$ são equivalentes.
  \item As formas $r(\cos\theta+i\sen\theta)$ e $re^{i\theta}$ são representações polares.
  \item A forma retangular favorece adição e subtração.
  \item A forma polar favorece multiplicação e divisão.
\end{itemize}
\begin{block}{Pergunta final}
Antes de calcular, qual representação torna a operação mais simples?
\end{block}
\end{frame}

\begin{frame}{Dever de casa}
\vspace{-0.55em}
\scriptsize
\begin{enumerate}
  \setlength{\itemsep}{2.5pt}
  \item Para $z=-3+4i$, determine $\operatorname{Re}(z)$, $\operatorname{Im}(z)$, $|z|$, um argumento e $\overline z$. Represente $z$ e $\overline z$ no plano.
  \item Encontre $x,y\in\mathbb{R}$ tais que $(2x-y)+(x+3y)i=7-i$.
  \item Calcule na forma $a+bi$:
  \[
    (3-2i)+(1+5i),\quad (2+i)(4-3i),\quad \frac{5+i}{2-i}.
  \]
  \item Determine $i^{73}$, $i^{202}$ e $i^{1001}$, justificando pelo resto da divisão por $4$.
  \item Converta $z=-2-2\sqrt3\,i$ para as formas $r(\cos\theta+i\sen\theta)$ e $re^{i\theta}$.
  \item Dados $z_1=4e^{i\pi/6}$ e $z_2=2e^{-i\pi/3}$, calcule $z_1z_2$ e $z_1/z_2$ nas formas polar e retangular.
  \item Resolva $x^2+4x+13=0$ em $\mathbb{C}$ e represente as raízes no plano complexo.
\end{enumerate}
\end{frame}

\begin{frame}{Referências}
\small

\begin{thebibliography}{99}

\bibitem{iezzi}
IEZZI, Gelson.
\textit{Fundamentos de Matemática Elementar: complexos, polinômios e equações}.
v. 6. São Paulo: Atual/Saraiva.

\bibitem{dante}
DANTE, Luiz Roberto.
\textit{Matemática: contexto \& aplicações: ensino médio}.
São Paulo: Ática.

\bibitem{bncc}
BRASIL. Ministério da Educação.
\textit{Base Nacional Comum Curricular: Ensino Médio}.
Brasília, DF: MEC, 2018.

\bibitem{openstax}
OPENSTAX.
\textit{Precalculus 2e}.
Houston: OpenStax, 2021.

\end{thebibliography}

\end{frame}

\begin{frame}
\centering
\vfill
{\FonteDestaque\LARGE Obrigado!}

\vspace{1em}
{\large Números Complexos}
\vfill
\end{frame}

\end{document}